%Autores: Prof. Me. Phyllipe Lima
%         Daylon Ramos da Silva
%Contato: phyllipe@inatel.br / phyllipe_slf@yahoo.com.br
%         biblioteca.pesquisa@inatel.br
%Modelo para escrita de artigos científicos para o curso de Pós-Graduação em Industria 4.0 do INATEL - Instituto Nacional de Telecomunicações. 
%Este template LaTeX é uma adaptação do modelo doc desenvolvido pelos professores Carlos Ynoguti e Dayan Guimarães

%Se você é novo no latex, um bom lugar para começar é
%https://pt.overleaf.com/learn
\documentclass[10pt,twocolumn]{article} 
\input{configuracoes/includes.tex}

%Configuracao do cabecalho
%Fique atento as datas e remova a cor vermelha!
\pagestyle{fancyplain}{%
    \fancyhf{} % limpa os campos
    \chead{ \textcolor{black!50}{\textbf{
    \textcolor{red}{X} SEMINÁRIO DE INDÚSTRIA 4.0\\ 
    INSTITUTO NACIONAL DE TELECOMUNICAÇÕES – INATEL \\
    \textcolor{red}{MÊS DA APRESENTAÇÃO} DE 20\textcolor{red}{XX} 
    \vspace*{2\baselineskip}
    }}}
    \rhead{\textcolor{black!50}{\textbf{
    e-ISSN XXXX-XXXX \hspace*{-1pt}\\
    }}}
    \cfoot{\thepage}
}
%Fim do cabeçalho

\AtBeginBibliography{\small}

%Inicio do documento
\begin{document}

\input{configuracoes/config.tex}
\title{\Huge \bf Instrução Para os Autores}

\author{\large 
Phyllipe Lima \&
Daylon Ramos da Silva }

\maketitle

\noindent\blfootnote{Trabalho de Conclusão de Curso apresentado ao Instituto Nacional de Telecomunicações, como parte dos requisitos para a obtenção do Certificado de Pós-Graduação em Indústria 4.0. Orientador: Prof. (XXXXX). Trabalho aprovado em XX/20XX.}

\input{secoes/01-resumo.tex}
\input{secoes/02-introducao.tex}
\input{secoes/03-instrucoes-gerais.tex}
\input{secoes/04-conclusao.tex}

\printbibliography

\input{secoes/05-biografia.tex}
\end{document}

